\documentclass[11pt]{article}
\usepackage[a4paper,margin=25mm]{geometry}
\usepackage{amsmath,amssymb,amsthm}
\usepackage[hidelinks]{hyperref}
\emergencystretch=2em
\newtheorem{theorem}{Theorem}
\newtheorem{lemma}{Lemma}
\newcommand{\Q}{\mathbb Q}
\newcommand{\Z}{\mathbb Z}
\title{Four tetrahedral spanning trees versus a fourth derivative of 24\\
\large Disproof of Conjecture 00000004222}
\author{gaochengzhecpu}\date{}
\begin{document}\maketitle
\begin{abstract}
The boundary of a tetrahedron has four two-dimensional spanning trees.
Its top-dimensional Laplacian acts on four faces, so its monic
characteristic polynomial has fourth derivative $4!=24$ everywhere,
including at $1$. This disproves the stated equality.
\end{abstract}

\section*{Statement and conventions}
The conjecture asserts that a higher-dimensional tree count equals the
derivative at $1$ of a Laplacian characteristic polynomial, with derivative
order equal to the number of faces in that dimension. We use the ordinary
derivative of the characteristic polynomial $\det(tI-L)$, and the
unweighted combinatorial Laplacian on that dimension's chains.
No rescaling or reduced characteristic polynomial is specified in the source.

For a two-dimensional complex with complete one-skeleton, a spanning
two-tree keeps all vertices and edges and selects triangular faces such
that the selected boundary map is injective over $\Q$ and its image is
$\ker\partial_1$. These are the chain descriptions of vanishing rational
$H_2$ and $H_1$. This agrees with the usual simplicial spanning-tree
conditions here; see Duval--Klivans--Martin, Section 1.2.\footnote{\href{https://www.dam.brown.edu/people/cklivans/SST.pdf}{A. M. Duval, C. J. Klivans, and J. L. Martin, \emph{Simplicial matrix-tree theorems}, Section 1.2.}}
We also check below that every tree in this example has trivial integral
$H_1$, so a squared torsion-order weight does not change the count.

\section*{The actual chain complex}
Let $K=\partial[0,1,2,3]$, the family of all proper subsets of four
vertices. Its six increasing oriented edges and four increasing oriented
faces are, respectively,
\[
 (01,02,03,12,13,23),\qquad (123,023,013,012).
\]
Use the usual boundaries $\partial_1[u,v]=[v]-[u]$ and
$\partial_2[a,b,c]=[b,c]-[a,c]+[a,b]$. Their matrices are
\[
 B_1=\begin{pmatrix}
 -1&-1&-1&0&0&0\\
 1&0&0&-1&-1&0\\
 0&1&0&1&0&-1\\
 0&0&1&0&1&1
 \end{pmatrix},\qquad
 B_2=\begin{pmatrix}
 0&0&1&1\\0&1&0&-1\\0&-1&-1&0\\
 1&0&0&1\\-1&0&1&0\\1&1&0&0
 \end{pmatrix}.
\]
Direct calculation gives $B_1B_2=0$. Set
$\sigma=(1,-1,1,-1)^T$. Solving $B_2x=0$ gives
\[
 x_2+x_3=0,\quad x_1-x_3=0,\quad x_0+x_3=0,
 \qquad x=x_0\sigma.
\]
Thus $\ker B_2=\Q\sigma$, and every nonzero two-cycle uses all four faces.
For any one-cycle $z=(z_0,\ldots,z_5)^T$, the first three rows of
$B_1z=0$ imply
\[
 z_2=-z_0-z_1,\qquad z_4=z_0-z_3,\qquad z_5=z_1+z_3.
\]
Consequently $a=(z_3,z_1,z_0,0)^T$ satisfies $B_2a=z$.

\begin{lemma}
A subset of the four triangular faces is a spanning two-tree precisely
when it omits exactly one face. Therefore $\tau_2(K)=4$.
\end{lemma}
\begin{proof}
If face $v$ is omitted, every supported two-cycle is a scalar multiple
of $\sigma$ with its $v$-coordinate zero, hence is zero.
For a one-cycle $z$, take the filling $a$ above and define
\[
 x=a-(a_v\sigma_v)\sigma.
\]
Since $\sigma_v^2=1$, this has $x_v=0$, while
$B_2x=B_2a=z$. Thus omitting one face gives a spanning two-tree.

Selecting every face fails injectivity because $\sigma\ne0$ is a cycle.
If distinct faces $v,w$ are omitted, let $e_v$ denote the $v$th coordinate
vector. The one-cycle $B_2e_v$ cannot have a supported filling $x$:
otherwise $B_2(x-e_v)=0$, so $x-e_v=c\sigma$.
Its $w$-coordinate forces $c=0$, whereas its $v$-coordinate is $-1$,
a contradiction. This proves the classification and count.
\end{proof}

All displayed filling formulas have integer coefficients. For integral
$z$, both $a$ and $x$ are integral, so every integral one-cycle is a
boundary in each of the four trees. The same kernel argument rules out
integral two-cycles. In particular, $H_1(T;\Z)=0$ for every counted tree
$T$, and the usual weight $|\operatorname{Tor}H_1(T;\Z)|^2$ equals $1$.
Geometrically, these four complexes are also the four closed disks
obtained by removing one open face of the tetrahedral sphere.

\section*{The prescribed derivative}
There are no three-dimensional faces, so the top-dimensional Hodge
Laplacian is exactly
\[
 L_2=B_2^TB_2\in M_4(\Q),\qquad p(t)=\det(tI_4-L_2).
\]
Every $4\times4$ matrix has a monic characteristic polynomial of degree
$4$. Hence $p(t)=t^4+c_3t^3+c_2t^2+c_1t+c_0$, and
\[
 \boxed{\tau_2(K)=4\ne24=p^{(4)}(1),\qquad f_2(K)=4.}
\]
One also commonly uses the up-down Laplacian $L_1^{\mathrm{up}}=B_2B_2^T$
on edges, as in the cited simplicial matrix-tree theorem. This convention
also contradicts the claimed derivative order. Indeed, multiplication
of the displayed matrix gives
\[
 B_2^TB_2=4I_4-\sigma\sigma^T,\qquad \sigma^T\sigma=4.
\]
Thus its eigenvalues are $4,4,4,0$. The $6\times6$ matrix $B_2B_2^T$
has the same nonzero eigenvalues and three zero eigenvalues. Therefore
\[
 p_{\mathrm{up}}(t)=t^3(t-4)^3=t^6-12t^5+48t^4-64t^3,
 \qquad p_{\mathrm{up}}^{(4)}(1)=360-1440+1152=72\ne4.
\]
Here the derivative order remains the number of two-faces, namely $4$.
This supplementary convention check is a paper calculation; the Lean
proof below formalizes the $4\times4$ top-dimensional Laplacian.

The contradiction concerns the derivative and derivative order stated in
the source. It does not invoke, or refute, a distinct identity involving
reduced Laplacians, normalized coefficients, or products of nonzero
eigenvalues. Even division by $4!$ would give $1$, not the count $4$.

\section*{Correspondence with Lean}
The finite family \texttt{tetrahedron} is proved downward closed.
The actual two-faces are enumerated bijectively by \texttt{faceSet};
the six increasing edges are enumerated by \texttt{tail} and \texttt{head}.
The matrices \texttt{B\_\allowbreak{}1} and \texttt{B\_\allowbreak{}2} (written with Unicode
subscripts in the source) are defined from the displayed simplicial
boundary formulas, before their entries are computed.
Their linear maps satisfy \texttt{boundary\_\allowbreak{}boundary}.

\texttt{IsSpanningTwoTree} uses those actual rational boundary maps and
the support of a selected-face chain. The theorem
\texttt{tree\_\allowbreak{}iff\_\allowbreak{}one\_\allowbreak{}missing} proves the classification; the finite
cardinality \texttt{treeCount} is then proved equal to $4$.
The object \texttt{characteristic} is Mathlib's actual
\texttt{Matrix.charpoly} of \texttt{L\_\allowbreak{}2}.
\texttt{fourth\_\allowbreak{}derivative} uses monicity, the polynomial degree, and
the coefficient formula for iterated derivatives. The final theorem
\texttt{tree\_\allowbreak{}count\_\allowbreak{}ne\_\allowbreak{}prescribed\_\allowbreak{}derivative} substitutes the actual
cardinality of the face type as the derivative order.
The formal proof is over $\Q$; the integral-homology and topological
interpretations above are supplied by the explicit integer filling
argument on paper, rather than claimed as imported Lean homology results.

\section*{Reproduction}
The project pins Lean 4.19.0 and Mathlib commit
\\\texttt{c44e0c8ee63ca166450922a373c7409c5d26b00b}.\par
From \texttt{lean/}, run \texttt{lake build} and
\texttt{lake env lean -DwarningAsError=true Main.lean}.
On a new machine, \texttt{lake exe cache get} retrieves the official
pinned dependency artifacts. The submission is checked in a fresh
build directory. Printed dependencies use only the standard
\texttt{propext}, \texttt{Classical.choice}, and \texttt{Quot.sound}.
No additional axioms, proof placeholders, or numerical oracles are used.
Run \texttt{tectonic main.tex} to reproduce the PDF; no auxiliary numerical
script is needed. Logs, hashes, and visual-review records are included
under \texttt{verification/}.

\section*{Conjecture source}
\href{https://github.com/The-Last-Math-Competition/The-Last-Math-Competition/blob/main/conjectures/00000004222.md}{The Last Math Competition, Conjecture 00000004222}.
The exact bilingual source and its commit and SHA-256 are preserved in
this submission.
\end{document}
